\section{Details in Overview}\label{sec:tech:details-overview}

We introduce how to use \tyName to specify the well-typedness of the De Bruijn STLC example introduced in the overview section.

\begin{figure}[th!]
  \begin{minted}[xleftmargin=5pt, numbersep=4pt, linenos = true,  escapeinside=??]{OCaml}
    type ty = Int | Arr of ty * ty
    type serialized_term = Const of int | Var of int | Abs of ty | EndAbs
                         | App | AppL | AppR
  \end{minted}
  \vspace*{-.1in}
  \caption{STLC terms in a serialized form.}
  \label{fig:stlc-syntax}
  \vspace*{-.1in}
\end{figure}

We explicitly define the syntax of the serialized STLC terms in \autoref{fig:stlc-syntax}. The definitions of constants and variables are standard, and variables
carry their de Bruijn index. However, the usual function abstraction
constructor ``$\Code{Abs\;\mykw{of}\;ty * term}$'' is now represented
by two operations, ``$\Code{Abs\;\mykw{of}\;ty}$'' and
``$\Code{EndAbs}$'', where the function body is the term that is
generated between them. Similarly, the typical application constructor
``$\Code{App\;\mykw{of}\;term * term}$'' is now captured by three
operations: ``$\Code{App}$'', ``$\Code{AppL}$'', and ``$\Code{AppR}$''.
Here, ``$\Code{App}$'' begins the application, ``$\Code{AppL}$'' marks the beginning of
the argument, and ``$\Code{AppR}$'' signals the end of the application.

An example of a serialized STLC term is shown below, where we color
the corresponding parts of the term and the serialized term for better
readability.
\begin{align*}
  \text{\textcolor{gray}{STLC term:}} \quad & \textcolor{blue}{(\lambda \Int}. \textcolor{orange}{[0]} \textcolor{blue}{)}\; \textcolor{DeepGreen}{3} \\
  \text{\textcolor{gray}{serialized STLC term:}} \quad & \eff{app}; \textcolor{blue}{\zevent{abs}{\Int}}; \textcolor{orange}{\zevent{var}{0}}; \textcolor{blue}{\eff{endAbs}}; \eff{appL}; \textcolor{DeepGreen}{\zevent{const}{3}}; \eff{appR}
\end{align*}
\noindent
We denote de Bruijn indices with ``$[...]$'' to distinguish them from
constants. The interpreter accepts a
serialized lambda term and evaluates it to a value. In the example
above, the input sequence should be reduced as $(\lambda \Int. [0])\;3
\hookrightarrow 3$.

\begin{figure}[t!]
  {\footnotesize
    \begin{align*}
      \eff{const} :~ &  x{:}\Int\sarr \uhat{\allA}{\Unit}{\msgA{const}{x}\msgAGray{ty}{\CInt}} \\
      \eff{appR} :~ & \uhat{\allA\msgC{app}\allA  \msgAGray{ty}{\Code{Arr}(t_1, t_2)}  \msgC{appL}  \allA  \msgAGray{ty}{t_1}  }{\Unit}{\msgC{appR}  \msgAGray{ty}{t_2}} \\
      \eff{endAbs} :~ & \uhat{\allA\msgAGray{depth}{d}\msgA{abs}{t_2}\msgAGray{depth}{d+1}\allA  \msgAGray{ty}{t_1}}{\Unit}{ \msgC{endAbs}\msgAGray{depth}{d}
        \msgAGray{putTy}{\Code{Arr}(t_1,t_2)}} \\
      \eff{var} :~ &  i{:}\urt{\Int}{\vnu \geq 0} \sarr \uhat{\underbrace{\allA\msgA{abs}{t}\msgAGray{depth}{d}(\anyA\setminus\msgC{endAbs})^*}_\text{type $t$ at the $d$ level is open}\underbrace{\msgAGray{depth}{d + i}}_\text{current level is $d + i$} }{\Unit}{\msgA{var}{i}\msgAGray{putTy}{t}  \allA}
    \end{align*}
  }
  \caption{\tyName{} specifications of STLC operations. } %
  \label{fig:stlc-spec}
\end{figure}

The types for STLC operations are shown in \autoref{fig:stlc-spec}. The type for $\eff{const}$ has no constraint on its prefix trace
($\allA$), but records a $\textcolor{gray}{\eff{ty}}$ event in
the trace that records the type of the output of the $\eff{const}$
expression as $\Code{Int}$. The arguments to a ghost event are either
serialized terms (like $\Code{Int}$) or ghost variables.
Type-checking a trace using ghost events involves solving constraints
over their arguments based on algebraic axiomatizations, examples of
which are given below. The history and future regex of the $\eff{appR}$ type in
\autoref{fig:stlc-spec} characterizes both the syntax and well-typing
constraints of an STLC function application. It requires that (a)
prior to evaluating a function argument ($\eff{appL}$), the type of
the abstraction is \Code{Arr(t$_1$,t$_2$)}, i.e., the abstraction
being applied is a function type ($\msgAGray{ty}{\Code{Arr}(t_1,t_2)}$); (b) after the argument is evaluated, its type is
  determined to be \Code{t$_1$} ($\msgAGray{ty}{t_1}$); and (c)
  the type of the application's result is \Code{t$_2$}
  ($\msgAGray{ty}{t_2}$).

The de Bruijn index constraint is captured via the $\eff{\textcolor{gray}{depth}}$ operation in the \tyName{}s of the $\eff{endAbs}$ and $\eff{var}$ operations. In addition to the well-typing constraint, the \tyName{} of the $\eff{endAbs}$ operation maintains the depth of nested function abstractions using $\eff{\textcolor{gray}{depth}}$. Specifically, its history regex obtains the current depth ($\msgAGray{depth}{d}$), increases it by one at the beginning of the function body, and restores it after the $\eff{endAbs}$ event.
On the other hand, a valid index $i$ in $\eff{var}$ must be derived from an open abstraction at some depth in the prefix trace. The history regex of the $\eff{var}$ operation precisely characterizes this constraint: the abstraction at depth $d$ must be open ($\allA\msgA{abs}{t}\msgAGray{depth}{d}(\anyA\setminus\msgC{endAbs})^*$), and the current depth must be $d + i$.

\section{Operational Semantics}\label{sec:tech:semantics}

The small-step operational semantics of our core language are shown in
\autoref{fig:semantics}.

\begin{figure}[t!]
    {\footnotesize
    \vspace*{-.15in}
    {\small\begin{flalign*}
      &\text{\textbf{Handler Semantics }}\ 
      \fbox{$\alpha \vDash \zevent{op}{\overline{c}} \Downarrow c$ \quad $\primop(\overline{c}) \Downarrow c$} &
      \text{\textbf{Operational Semantics }}\ 
      \fbox{$\steptr{\alpha}{e}{\alpha}{e}$}
    \end{flalign*}}
        \begin{prooftree}
        \mhypo{30mm}
        {
          $\alpha \vDash \zevent{op}{\overline{c}} \Downarrow c'$
        }
          \infer1[\textsc{\small StEfOp}]{
            \steptr{\alpha}{\eff{op}\;\overline{c}}{[\zevent{op}{\overline{c}\;c'}]}{c'}
          }
        \end{prooftree}
          \quad
          \mprooftr{35mm}{
            $\primop\ \overline{c} \Downarrow c'$
            }{StOp}
            {$\steptr{\alpha}{\primop\ \overline{c}}{[]}{c'}$}
            \quad
    \\[0.8em]
    \mprooftr{25mm}{
    $\phi\Downarrow \top$
    }{StAssume}{
    $\steptr{\alpha}{\sassume{\phi}}{[]}{()}$
    }
    \quad
    \mprooftr{25mm}{$\phi \Downarrow \top$}{StAssert}{
    $\steptr{\alpha}{\sassert{\phi}}{[]}{()}$
    }
    \quad
    \mprooftr{35mm}{$\phi \Downarrow \bot$}{StAssert}{
    $\steptr{\alpha}{\sassert{\phi}}{[]}{\Code{raise\;Error}}$
    }
    \\[.8em]
    \mprooftr{65mm}{
      $\steptr{\alpha}{e_1}{\alpha'}{e_1'}$
    }{StLetE1}{
    $\steptr{\alpha}{\zlet{y}{e_1}{e_2}}{\alpha'}{\zlet{y}{e_1'}{e_2}}$
    }\quad
    \mprooftr{35mm}{}{StChoice1}{
    $\steptr{\alpha}{(e_1 \oplus e_2)}{[]}{(e_1)}$
    }
    \\[0.8em]
    \mprooftr{55mm}{}{StLetE2}{
    $\steptr{\alpha}{\zlet{y}{v}{e_2}}{[]}{e_2[y\mapsto v]}$
    }
    \quad
    \mprooftr{35mm}{}{StChoice2}{
    $\steptr{\alpha}{(e_1 \oplus e_2)}{[]}{(e_2)}$
    }
    \\[0.8em]
    \mprooftr{40mm}{}{StAppLam}{
    $\steptr{\alpha}{\zlam{x}{s}{e_1}\ v_x}{[]}{e_1[x\mapsto v_x]}{e_2}$
    }
    \\[0.8em]
    \mprooftr{80mm}{}{StAppFix}{
    $\steptr{\alpha}{(\zfix{f}{s}{x}{s_x}{e_1})\ v_x}{[]}{(\zlam{f}{s}{e_1[x\mapsto v_x]}) \ (\zfix{f}{s}{x}{s_x}{e_1})}$
    }
}
\caption{Full Operational Semantics}
\label{fig:semantics}
\end{figure}

\newpage
\section{Basic Typing Rules}\label{sec:tech:basic-typing}

The basic typing rules of our core language and qualifiers are shown
in \autoref{fig:basic-type-rules}. We use an auxiliary function
$\Code{Ty}$ to provide a basic type for the primitives of our language,
e.g., constants, built-in operators, and data constructors.

\begin{figure}[h!]
 {\footnotesize
 {\small
\begin{flalign*}
 &\text{\textbf{Basic Typing }} & \fbox{$\Gamma \basicvdash e : s$ \quad $\Gamma \basicvdash \eff{op} : s$ \quad $\Gamma \basicvdash \primop : s$}
\end{flalign*}}
\\ \
\begin{prooftree}
\hypo{}
\infer1[\textsc{BtConst}]{
\Gamma \basicvdash c : \Code{Ty}(c)
}
\end{prooftree}
\quad
\begin{prooftree}
\hypo{ \Gamma(x) = s }
\infer1[\textsc{BtVar}]{
\Gamma \basicvdash x : s
}
\end{prooftree}
\quad
\begin{prooftree}
\hypo{ \Gamma, x{:}b \basicvdash e : s }
\infer1[\textsc{BtFun}]{
\Gamma \basicvdash \zlam{x}{b}{e} : b\sarr s
}
\end{prooftree}
\quad
\begin{prooftree}
  \hypo{ \Gamma, f{:}b\sarr s, x{:}b \basicvdash e : s }
  \infer1[\textsc{BtFix}]{
  \Gamma \basicvdash \zfix{f}{b\sarr s}{x}{b}{e} : b\sarr s
  }
\end{prooftree}
\\[0.8em]
\begin{prooftree}
  \hypo{\Code{Ty}(\eff{op}) = s}
  \infer1[\textsc{BtEfOp}]{
  \Gamma \basicvdash \eff{op} : s
  }
  \end{prooftree}
\quad
\begin{prooftree}
  \hypo{\Code{Ty}(\primop) = s}
  \infer1[\textsc{BtPrimOp}]{
  \Gamma \basicvdash \primop : s
  }
\end{prooftree}
\quad
\begin{prooftree}
  \hypo{\Gamma \basicvdash e_1 : s
  \quad \Gamma \basicvdash e_2 : s}
  \infer1[\textsc{BtChoice}]{
  \Gamma \basicvdash e_1 \oplus e_2 : s
  }
\end{prooftree}
\\[0.8em]
\begin{prooftree}
  \hypo{\Gamma \basicvdash \eff{op}: \overline{b_i}\sarr b \quad \forall i.\Gamma \basicvdash c_i : b_i}
  \infer1[\textsc{BtEfOpApp}]{
  \Gamma \basicvdash \eff{op}\;\overline{c_i} : b
  }
  \end{prooftree}
\quad
\begin{prooftree}
\hypo{\Gamma \basicvdash \primop: \overline{s_i}\sarr s \quad \forall i.\Gamma \basicvdash v_i : s_i}
\infer1[\textsc{BtPrimOpApp}]{
\Gamma \basicvdash \primop\,\overline{v_i} : s
}
\end{prooftree}
\\[0.8em]
\begin{prooftree}
\hypo{\Gamma \basicvdash \phi : \Code{bool}}
\infer1[\textsc{BtAssume}]{
\Gamma \basicvdash \sassume{\phi} : \Unit
}
\end{prooftree}
\quad
\begin{prooftree}
\hypo{\Gamma \basicvdash \phi : \Code{bool}}
\infer1[\textsc{BtAssert}]{
\Gamma \basicvdash \sassert{\phi} : \Unit
}
\end{prooftree}
\\[0.8em]
\begin{prooftree}
  \hypo{\Gamma\basicvdash e : s \quad \Gamma,y{:}s \basicvdash e' : s}
  \infer1[\textsc{BtLetE}]{
  \Gamma \basicvdash \zlet{y}{e}{e'} : s
  }
\end{prooftree}
  \quad
\begin{prooftree}
  \hypo{\Gamma \basicvdash v_1 : s\sarr s' \quad \Gamma \basicvdash v_2 : s}
  \infer1[\textsc{BtApp}]{
  \Gamma \basicvdash v_1\;v_2 : \Unit
  }
\end{prooftree}
}
{\footnotesize
 {\small
\begin{flalign*}
 &\text{\textbf{Basic Qualifier Typing }} & \fbox{$\Gamma \basicvdash \phi: s$}
\end{flalign*}}
\\ \
\begin{prooftree}
\hypo{\Code{Ty}(c) = s }
\infer1[\textsc{BtLitConst}]{
\Gamma \basicvdash c : s
}
\end{prooftree}
\quad
\begin{prooftree}
\hypo{ \Gamma(x) = s }
\infer1[\textsc{BtLitVar}]{
\Gamma \basicvdash x : s
}
\end{prooftree}
\quad
\begin{prooftree}
\hypo{ }
\infer1[\textsc{BtTop}]{
\Gamma \basicvdash \top : \Code{bool}
}
\end{prooftree}
\quad
\begin{prooftree}
\hypo{ }
\infer1[\textsc{BtBot}]{
\Gamma \basicvdash \bot : \Code{bool}
}
\end{prooftree}
\\ \ \\ \ \\
\begin{prooftree}
\hypo{\Code{Ty}(\primop) = \overline{s_i}{\sarr}s \quad \forall i. \Gamma \basicvdash l_i : s_i }
\infer1[\textsc{BtLitOp}]{
\Gamma \basicvdash \primop\,\overline{l_i} : s
}
\end{prooftree}
\quad
\begin{prooftree}
\hypo{
\Gamma \basicvdash \phi : \Code{bool} }
\infer1[\textsc{BtNeg}]{
\Gamma \basicvdash \neg \phi : \Code{bool}
}
\end{prooftree}
\\ \ \\ \ \\
\begin{prooftree}
\hypo{
\Gamma \basicvdash \phi_1 : \Code{bool} \quad
\Gamma \basicvdash \phi_2 : \Code{bool} }
\infer1[\textsc{BtAnd}]{
\Gamma \basicvdash \phi_1 \land \phi_2 : \Code{bool}
}
\end{prooftree}
\quad
\begin{prooftree}
\hypo{
\Gamma \basicvdash \phi_1 : \Code{bool} \quad
\Gamma \basicvdash \phi_2 : \Code{bool} }
\infer1[\textsc{BtOr}]{
\Gamma \basicvdash \phi_1 \lor \phi_2 : \Code{bool}
}
\end{prooftree}
\quad
\begin{prooftree}
\hypo{
\Gamma, x{:}b \basicvdash \phi : \Code{bool} }
\infer1[\textsc{BtForall}]{
\Gamma \basicvdash \forall x{:}b. \phi : \Code{bool}
}
\end{prooftree}
}
\caption{Basic Typing Rules and Qualifier Typing Rules}
    \label{fig:basic-type-rules}
\end{figure}
